\documentclass[conference]{IEEEtran}
\IEEEoverridecommandlockouts
\usepackage{cite}
\usepackage{amsmath,amssymb}
\usepackage{booktabs}
\usepackage{graphicx}
\usepackage{url}

\begin{document}

\title{model2rtl: Compiling a Quantized Neural Network into\\ Portable RTL with a Shared Multiply-Select-Add Fabric}

\author{\IEEEauthorblockN{Rithwik}
\IEEEauthorblockA{\textit{Independent Research} \\
Email: sniyaantony@gmail.com}}

\maketitle

\begin{abstract}
Deploying a neural network as dedicated hardware normally means instantiating
one multiplier per synapse, or accepting the area and control cost of a general
matrix engine. We describe \emph{model2rtl}, a compiler that emits portable,
synthesizable Verilog-2001 for a quantized multilayer perceptron using a
\emph{Multiply-Select-Add} (MSA) fabric. Because every weight is one of
$K{=}16$ fixed integer levels, an activation $x_i$ can participate in only 16
distinct products regardless of its fanout. The fabric computes those 16
products once per cycle and lets every output neuron select the one its 4-bit
weight index names. Execution is input-serial and output-parallel, so a single
16-entry product bank serves both layers of the network.
The design separates the compute fabric from the trained parameters
completely: regenerating the fabric from a different weight set yields a
byte-identical file, which we verify on every compilation. We validate a
784--32--10 MNIST network end to end. Quantization to 4-bit weight indices and
uint8 activations costs 0.07 percentage points of test accuracy
(96.52\%~$\rightarrow$~96.45\%). Over 500 test images the RTL reproduces a
pure-integer reference model with zero mismatches across 16{,}000 hidden
activations and 5{,}000 logits, supported by 178{,}840 cycle-level internal
checks. The \emph{same} source then synthesizes through an FPGA-oriented and a
generic/ASIC-oriented Yosys flow with no source change, and both gate-level
netlists reproduce the integer model exactly on the same 500 images. Neither
netlist retains a single multiplier or DSP cell. We additionally generate seven
SKY130 mask-ROM macros for the parameters and verify all 102{,}640 programmed
cells against the canonical images, and we report a storage-area sweep in which
the hard ROM does not beat synthesized standard-cell storage anywhere in the
measured 1{,}024--50{,}176-bit range.
\end{abstract}

\begin{IEEEkeywords}
neural network accelerators, quantization, RTL generation, weight sharing,
multiplier-free inference, open-source EDA
\end{IEEEkeywords}

\section{Introduction}

A trained neural network is a fixed function. Once training stops, the weights
never change again, and in principle nothing about the inference datapath needs
to be general. This observation motivates a long line of work on turning
networks directly into circuits~\cite{finn,logicnets,hls4ml}.

The obstacle is arithmetic. A dense layer with $n$ inputs and $m$ outputs
implies $nm$ multiplications. Building them spatially is prohibitive for all
but the smallest networks; time-multiplexing them onto a smaller array
reintroduces the control, buffering and scheduling machinery the direct
approach was meant to avoid.

Aggressive weight quantization changes the shape of the problem. If every
weight is drawn from a small alphabet of $K$ integer levels, then for a given
activation $x_i$ there are only $K$ distinct products $x_i \cdot a_k$ that can
ever be required --- no matter how many neurons $x_i$ feeds. The multiplication
is a property of the \emph{activation}, not of the synapse. Each synapse then
needs only to \emph{select} among precomputed products, which costs a
multiplexer rather than a multiplier. We refer to this arrangement as
Multiply-Select-Add, a term that also appears in public trade coverage of
commercial hardwired-inference silicon~\cite{taalas}.

This paper contributes:

\begin{enumerate}
\item An RTL generator for a Multiply-Select-Add fabric that is
\emph{provably independent of the trained parameters}: substituting a different
weight and bias set yields a byte-identical Verilog file, checked by SHA-256 on
every compilation.
\item An integer arithmetic contract fixed analytically \emph{before} any RTL
exists, containing no multiplicative requantization scale, which is what makes
that independence achievable.
\item A verification methodology in which a pure-NumPy integer model is the
sole oracle at every level --- behavioral simulation, post-synthesis gate-level
simulation on two independent targets, and physical mask-ROM contents.
\item Measured results, including a negative one: the hard ROM we generate is
larger than synthesized standard-cell storage at every size we measured.
\end{enumerate}

We are explicit about scope. This work produces and verifies RTL. It does not
perform place-and-route, timing analysis or a full physical implementation, and
it makes no claim about clock frequency or finished-chip area.

\section{Related Work}

\textbf{Quantization.} Jacob et al.~\cite{jacob} established integer-only
inference with a multiplicative rescale between layers, which underpins TFLite
quantization. Our contract deliberately omits that multiplier: the only
requantization operator is a fixed power-of-two shift. This is a restriction,
but it is precisely what allows the fabric to contain no trained value, since a
learned scale would otherwise have to be embedded in the datapath.
Han et al.~\cite{deepcompression} introduced weight sharing via $k$-means
clustering, reducing AlexNet storage $35\times$; their shared weights are a
storage-side compression, whereas we exploit the same redundancy on the
\emph{compute} side.

\textbf{Computational reuse.} UCNN~\cite{ucnn} is the closest prior work in
spirit: it exploits weight repetition to reuse dot-product sub-computations,
generalizing sparsity (repeated zeros) to repeated values, and reports
$1.2\times$--$4\times$ energy improvement over an Eyeriss-style baseline. UCNN
discovers repetition dynamically in a general accelerator. We instead
\emph{fix} the alphabet at $K{=}16$ by construction, which makes the reuse
structure static and lets it be baked into the RTL rather than detected at run
time.

\textbf{Multiplier-free inference.} DeepShift~\cite{deepshift} constrains
weights to powers of two so multiplication becomes shift-and-sign-flip.
Our alphabet $\{-8,\dots,+7\}$ is not restricted to powers of two, yet we
observe the same outcome empirically: synthesis removes every multiplier
anyway, because one operand is a small constant (Section~\ref{sec:constmul}).

\textbf{LUT-native and circuit-level networks.} LUTNet~\cite{lutnet} treats the
FPGA $K$-LUT itself as the inference operator, achieving roughly twice the area
efficiency of binarized baselines. LogicNets~\cite{logicnets} goes further,
training quantized sparse networks whose neurons are directly equivalent to
truth tables and emitting a netlist of LUTs. Both are strongly FPGA-shaped: the
artifact is a LUT netlist. Our output is vendor-neutral Verilog-2001, which we
demonstrate synthesizes unchanged through both an FPGA flow and a generic
gate-level flow.

\textbf{Toolflows.} FINN~\cite{finn} builds streaming FPGA dataflow
accelerators for binarized networks, reporting 12.3M MNIST classifications/s on
a ZC706. hls4ml~\cite{hls4ml} compiles small networks to FPGA firmware via HLS
for sub-microsecond physics triggers. Both target FPGAs and rely on
vendor-specific backends; both are considerably more capable than this work in
network coverage. Our contribution is orthogonal: a single RTL source
demonstrated portable across FPGA-oriented and ASIC-oriented synthesis, with
post-synthesis functional equivalence verified on both.

\begin{table}[t]
\centering
\caption{Where each approach takes the saving. All target the same bottleneck:
inference arithmetic dominates area and energy.}
\label{tab:related}
\footnotesize
\begin{tabular}{@{}llp{2.05cm}@{}}
\toprule
Work & Lever & Artifact \\
\midrule
Deep Compression~\cite{deepcompression} & shared weights & compressed model \\
Jacob et al.~\cite{jacob} & integer-only math & quantized graph \\
DeepShift~\cite{deepshift} & powers of two & shift kernels \\
UCNN~\cite{ucnn} & run-time repetition & ASIC accelerator \\
LUTNet~\cite{lutnet} & FPGA LUT as operator & LUT netlist \\
LogicNets~\cite{logicnets} & neuron $\equiv$ truth table & LUT netlist \\
FINN~\cite{finn} & binarized dataflow & FPGA bitstream \\
hls4ml~\cite{hls4ml} & HLS from model & FPGA firmware \\
\midrule
\textbf{This work} & \textbf{static $K$-level bank} & \textbf{portable Verilog} \\
\bottomrule
\end{tabular}
\end{table}

\textbf{Synthesis across these works} shares a common motivation --- inference
arithmetic is the bottleneck and quantization is the lever. They diverge on
\emph{where} the saving is taken: in storage~\cite{deepcompression}, in the
operator~\cite{deepshift}, in the target's native primitive~\cite{lutnet,
logicnets}, or in run-time reuse~\cite{ucnn}. The strongest results are
FPGA-specific. The gap we address is portability with evidence: to our
knowledge, prior work does not demonstrate one unmodified RTL source passing
gate-level functional verification after two independent synthesis flows.

\section{Arithmetic Contract}

The integer specification was fixed analytically before any RTL was written and
never changed thereafter. Weight indices are 4 bits; the alphabet is the plain
two's-complement range $a_k = k - 8$, $k \in [0,15]$, giving levels $-8$ to
$+7$. Activations are unsigned 8-bit with zero point 0.

For layer 1, with $\mathbf{x} \in [0,255]^{784}$:
\begin{equation}
\mathrm{acc}_1[j] = \sum_{i} x_i \cdot a_{W_1[i][j]} + b_1[j]
\end{equation}
\begin{equation}
h[j] = \mathrm{clamp}\!\left(\frac{\max(\mathrm{acc}_1[j],0) + 128}{256},\, 0,\, 255\right)
\label{eq:requant}
\end{equation}
where the division is a round-half-up arithmetic right shift by 8. Layer 2
stops after the bias; the prediction is $\arg\max$ over signed logits, with
ties broken by lowest index.

Widths are derived from the alphabet and layer sizes rather than assumed:
products are 12-bit signed, layer-1 dot products 22-bit, layer-1 accumulators
23-bit, layer-2 dot products 17-bit and accumulators 18-bit.

The critical property of (\ref{eq:requant}) is what it \emph{lacks}. There is no
multiplicative requantization scale anywhere in the datapath. The shift is an
architectural constant, not a learned one. Consequently the fabric depends only
on topology, $K$, the activation width and the shift --- never on a trained
value.

Quantization-aware training with straight-through estimators simulating the
exact integer pipeline yields 96.45\% MNIST test accuracy against a 96.52\%
float baseline, a loss of 0.07 percentage points. All 16 alphabet levels are
used in both layers. Weight saturation during training is 0.92\% in layer~1 and
7.5\% in layer~2; the higher layer-2 figure reflects its 320 synapses and wider
per-synapse dynamic range and does not degrade final accuracy.

\section{The Multiply-Select-Add Fabric}

\begin{figure}[t]
\centering
\footnotesize
\begin{verbatim}
              activation x_i (uint8)
                        |
    +---------+---------+---------+---------+
    |         |         |         |         |
 x_i*-8    x_i*-7     .....    x_i*+6    x_i*+7
    |         |         |         |         |
    +------ 16 shared products (12b) -------+
                        |
       +--------+-------+--------+--------+
       |        |                |        |
    16:1 mux 16:1 mux   .....  16:1 mux   <- 4-bit
       |        |                |           weight
     acc[0]   acc[1]           acc[N]        index
\end{verbatim}
\caption{One product bank of $K{=}16$ generators is shared by every output
neuron of the active layer, and reused across input cycles and across both
layers. Each synapse contributes a multiplexer, not a multiplier.}
\label{fig:msa}
\end{figure}

Fig.~\ref{fig:msa} shows the datapath. Execution is \emph{input-serial,
output-parallel}: one activation enters per cycle, the 16 products are formed
combinationally, and all neurons of the active layer accumulate in parallel.

Three operation counts must be distinguished, and conflating them is the most
common way to overstate this kind of architecture:

\begin{table}[h]
\centering
\caption{Source-level operation counts for 784--32--10. None is a physical
multiplier count.}
\label{tab:counts}
\begin{tabular}{lr}
\toprule
Organization & Products \\
\midrule
Naive fully spatial (one per synapse) & 25{,}408 \\
Fully spatial MSA & 13{,}056 \\
\textbf{Implemented (input-serial MSA)} & \textbf{16} \\
\bottomrule
\end{tabular}
\end{table}

The implemented count is 16 because the bank is recomputed each cycle for the
current activation rather than unrolled across inputs. The cost is latency:
\begin{equation}
T = n_{\mathrm{in}} + 2n_{\mathrm{hid}} + n_{\mathrm{out}} + 6 = 864~\text{cycles}
\end{equation}
per inference, data-independent. Elaborating the generated Verilog confirms the
structure: exactly 16 \texttt{\$mul} cells and 42 selector instances appear in
the netlist, matching 32 layer-1 plus 10 layer-2 neurons.

Notably, sharing is not universally cheaper. It reduces product generators only
when a layer's output fanout exceeds $K$. Layer~1 (fanout 32) halves them;
layer~2 (fanout 10) would need $1.6\times$ \emph{more} in a fully spatial
arrangement. The input-serial organization sidesteps this because the same 16
generators serve both layers.

\subsection{When does sharing pay?}

Sharing reduces product generators only when a layer's output fanout exceeds
$K$. For a layer with $n$ inputs and $m$ outputs the naive spatial count is
$nm$ while the fully spatial MSA count is $nK$, so sharing wins iff $m > K$.
Layer~1 ($m{=}32 > 16$) halves its generators; layer~2 ($m{=}10 < 16$) would
need $1.6\times$ \emph{more}. The crossover sits exactly at fanout $= K$.

The input-serial organization removes this concern entirely, because the same
16 generators are reused across input cycles \emph{and} across both layers: the
implemented count is $K$ regardless of either layer's fanout. This is a
practical argument for time-multiplexing beyond area alone --- it makes the
architecture insensitive to layer shape.

The premise also requires that the alphabet actually be used.
Table~\ref{tab:hist} confirms it: all 16 levels are populated in both layers,
with the expected concentration near zero.

\begin{table}[t]
\centering
\caption{Weight-index distribution after quantization-aware training
(representative levels shown). All 16 levels are populated in both layers.}
\label{tab:hist}
\footnotesize
\begin{tabular}{@{}lrrrrrrrr@{}}
\toprule
Level & $-8$ & $-6$ & $-4$ & $-2$ & $0$ & $+2$ & $+4$ & $+7$ \\
\midrule
Layer 1 & 228 & 312 & 1013 & 2572 & 4201 & 2868 & 997 & 253 \\
Layer 2 & 23 & 8 & 18 & 16 & 22 & 28 & 26 & 8 \\
\bottomrule
\end{tabular}
\end{table}

\subsection{Weight independence}

The fabric is generated from topology alone. To verify this rather than assert
it, the compiler regenerates the fabric from a random parameter set of the same
shape on every run and compares SHA-256. For the reference model the fabric
hashes to \texttt{7757362642b37fd0\ldots} under the trained weights, under a
substituted weight set, and under substituted biases. Compilation fails closed
if these ever differ. Only the parameter ROM is model-dependent.

\section{Parameter Storage}

Parameters are packed into four canonical images: weights $784{\times}128$ and
$32{\times}40$ bits (eight 4-bit indices per 32-bit group, neuron 0 in the
least significant field), and biases $32{\times}22$ and $10{\times}17$ signed.
These images are hashed and are the single source of truth for every backend,
which makes it structurally impossible to physically build one dataset while
testing another.

Two backends implement the same fixed interface --- synchronous read, one-cycle
latency, enable-gated hold --- so the fabric cannot tell which is attached:
(i) a portable Verilog-2001 case/constant ROM usable on both FPGA and ASIC, and
(ii) SKY130 mask-ROM macros generated with OpenRAM's ROM compiler~\cite{openram}.

\section{Verification Methodology}

A pure-NumPy integer model implementing the contract of Section~III is the sole
oracle at every stage. Floating-point Keras output is a reference number only
and is never used to check RTL arithmetic. Verification proceeds at four
levels: behavioral RTL, post-synthesis gate level on two targets, physical ROM
contents, and full-model inference through the physical organization.

\section{Results}

\subsection{Behavioral verification}

Over the first 500 images of the official MNIST test set, taken in order with
no filtering, the RTL reproduces the integer model exactly: 0 mismatches across
16{,}000 hidden activations, 5{,}000 logits and 500 predictions, on both
backends and at 864 cycles every image. Cycle-level internal checkpointing over
20 images contributes 178{,}840 further comparisons with 0 failures, including
per-cycle proof that each consumed ROM word corresponds to the address issued
exactly one cycle earlier. Stall, reset, back-to-back, argmax-tie and
arithmetic-edge suites are all clean.

\begin{table}[t]
\centering
\caption{Verification levels. The same pure-integer model is the oracle
throughout.}
\label{tab:verif}
\footnotesize
\begin{tabular}{@{}lrr@{}}
\toprule
Level & Comparisons & Mismatches \\
\midrule
Behavioral RTL, 500 images & 21{,}500 & 0 \\
Cycle-level internal trace, 20 images & 178{,}840 & 0 \\
Post-synthesis GLS, iCE40 & 5{,}500 & 0 \\
Post-synthesis GLS, generic & 5{,}500 & 0 \\
Physical ROM programmed cells & 102{,}640 & 0 \\
Weight indices via physical banks & 25{,}408 & 0 \\
\bottomrule
\end{tabular}
\end{table}

Input timing was varied to confirm the handshake is not accidentally
order-dependent. With no stalls an inference takes 864 cycles; with a bubble
after every seventh accepted activation, 975 cycles; under a pseudo-random LFSR
pattern, 1{,}223--1{,}297 cycles. Results are bit-identical in all three cases
and only latency changes. Reset was injected at six points with no stale state
surviving, and 500 inferences ran back to back in one simulation with no reset
between them.

\subsection{Dual-target synthesis portability}

The same four RTL files, byte-identical, were synthesized through
\texttt{synth\_ice40} and through an independent generic flow
(\texttt{proc}/\texttt{flatten}/\texttt{opt}/\texttt{memory}/\texttt{techmap}/
\texttt{abc -g simple}) using Yosys~\cite{yosys}. No source was patched. Both
netlists were then gate-level simulated with the official cell models.

\begin{table}[h]
\centering
\caption{Post-synthesis gate-level results, 500 MNIST images per target.}
\label{tab:gls}
\begin{tabular}{lrr}
\toprule
& iCE40 & Generic \\
\midrule
Logits compared & 5{,}000 & 5{,}000 \\
\textbf{Logit mismatches} & \textbf{0} & \textbf{0} \\
\textbf{Prediction mismatches} & \textbf{0} & \textbf{0} \\
Cycles per inference & 864 & 864 \\
\midrule
Total cells & 9{,}079 & 45{,}707 \\
LUT4 / gates & 6{,}429 & 40{,}614 \\
Flip-flops & 1{,}614 & 1{,}742 \\
Carry cells & 1{,}004 & --- \\
Block RAMs & 32 & 0 \\
\textbf{Multiplier / DSP cells} & \textbf{0} & \textbf{0} \\
\bottomrule
\end{tabular}
\end{table}

Table~\ref{tab:gls} summarizes. The two netlists also agree with each other bit
for bit. Both flows are byte-deterministic across clean re-runs.

Two observations deserve emphasis. First, the portable case-ROM was
\emph{automatically} inferred into 32 iCE40 block RAMs with no FPGA-specific
RTL, vendor primitive or synthesis pragma. Second, the architectural latency of
864 cycles survived both flows: synthesis added no pipeline stage.

\subsection{What synthesis did to the multiplications}
\label{sec:constmul}

The source contains exactly 16 \texttt{*} operators. Both netlists contain
\emph{zero} multiplier or arithmetic cells. Inspecting the netlists rather than
assuming, of the 192 product-bank bits, 37 (FPGA) are literal constants and 41
are plain aliases of the activation register --- 41\% of the bank costs no
logic at all. Where the FPGA flow preserved wire names, the drivers are
explicit: $x{\cdot}0$ folded to a literal zero, $x{\cdot}1$, $x{\cdot}2$ and
$x{\cdot}4$ became pure wiring with constant zero fill, and $x{\cdot}(-8)$ is a
shift of a shared negated value. The remaining 114 bits have no driver at all:
synthesis fused product generation into the 16:1 selection.

The honest statement is therefore not that we reduced 25{,}408 multipliers to
16. It is that \emph{the architecture exposes only sixteen constant-weight
product alternatives per activation, and synthesis then eliminates explicit
multiplier hardware entirely}. The three counts in Table~\ref{tab:counts} are
source-level quantities; the synthesized results are a different kind of
measurement.

\subsection{Physical ROM generation}

Seven SKY130 macros were generated. Two physical transformations were needed
because the ROM compiler expresses word size in bytes and could not route a
$784{\times}128$ array: layer-1 weights were banked into four parallel
$784{\times}32$ macros sharing one address, and both bias memories were stored
as 24-bit sign-extended words. Both transformations are exactly reversible, and
the reverse map is an automated invariant.

Every programmed cell was read back out of the generated SPICE netlists and
compared against the canonical images: \textbf{102{,}640 / 102{,}640 exact}.
All 858 logical rows reconstruct from the banks and all 25{,}408 weight indices
survive. The physical organization has zero functional effect: over the same
500 images the physical backend gives 0 hidden, logit and prediction
mismatches at 864 cycles.

\subsection{Storage area}

\begin{table}[h]
\centering
\caption{Storage-size sweep. Identical deterministic contents; OpenROM area is
a GDS bounding box, portable area is a liberty cell-area sum. These are
different kinds of area.}
\label{tab:crossover}
\begin{tabular}{rrrr}
\toprule
Bits & OpenROM ($\mu$m$^2$) & Portable ($\mu$m$^2$) & Ratio \\
\midrule
1{,}024  & 13{,}275 & 1{,}964  & 6.76 \\
4{,}096  & 18{,}467 & 4{,}773  & 3.87 \\
16{,}384 & 38{,}879 & 13{,}753 & 2.83 \\
25{,}088 & 53{,}669 & 18{,}434 & 2.91 \\
50{,}176 & 92{,}011 & 30{,}935 & 2.97 \\
\bottomrule
\end{tabular}
\end{table}

The seven macros total 252{,}382~$\mu$m$^2$ of GDS bounding box. The same
parameters mapped to SKY130 standard cells occupy 9{,}406 cells and
58{,}336~$\mu$m$^2$ of liberty cell area, a raw ratio of $4.33\times$.

Table~\ref{tab:crossover} sweeps storage size. \textbf{No crossover was
measured}: the portable mapping is smaller at every point from 1{,}024 to
50{,}176 bits, and the ratio flattens near $2.9\times$ rather than converging
toward 1. We state this narrowly. It does not establish that a hard ROM never
wins; it establishes that no OpenROM area advantage was observed over the
measured range with this generator and library. As context only --- not a
measurement --- the two would break even at the deepest point if the portable
block were placed at 33.6\% utilization or worse.

\section{Generalization Beyond the Reference Network}

The architecture is not specific to 784--32--10. The compiler accepts any
two-layer dense network: it infers the topology from the model, derives all
datapath widths from the alphabet and layer sizes, and emits the corresponding
fabric. To test this on a model the project had never seen, we trained a
784--64--10 network on Fashion-MNIST and compiled it.

Post-training quantization alone is not free at 4 bits, because the float
weights were never trained to survive a 16-level alphabet. Fine-tuning with the
exact integer pipeline in the forward pass recovers the loss and more.
Table~\ref{tab:fashion} reports accuracy on a held-out split the quantizer never
saw.

\begin{table}[h]
\centering
\caption{Fashion-MNIST 784--64--10, held-out split. RTL was simulated, not
merely elaborated.}
\label{tab:fashion}
\footnotesize
\begin{tabular}{@{}lr@{}}
\toprule
& Accuracy \\
\midrule
Float Keras baseline & 83.0\% \\
Post-training quantization $\rightarrow$ RTL & 82.7\% \\
Quantization-aware fine-tuning $\rightarrow$ RTL & \textbf{86.0\%} \\
\midrule
RTL vs.\ integer model (3{,}000 logits) & \textbf{0 mismatches} \\
Cycles per inference & 928 \\
\bottomrule
\end{tabular}
\end{table}

Two details matter. The requantization shift is no longer an architectural
constant once the topology is free: it must be chosen per model. We choose it
by measurement, sweeping candidate shifts against calibration data. The sweep
exposes a clear trade --- at shift~4, 43\% of hidden activations saturate and
accuracy collapses to 74.8\%; at shift~12, 48\% are driven to zero and accuracy
falls to 96.2\%; the plateau sits at 7--9. Assuming a value here would silently
cost several points.

The compiler also refuses what it cannot build. A three-layer MLP and a
convolutional model are rejected with the layers they actually contain, rather
than partially compiled. In one downloaded case a three-layer network had
dimensionally compatible first and last layers, so dropping the middle layer
would have produced hardware that elaborates cleanly and computes a different
network. Refusing is the correct behavior.

\section{Discussion}

The enabling decision in this work is arithmetic, not architectural. Fixing the
requantization operator to a power-of-two shift, and admitting no multiplicative
scale, is what makes the compute fabric provably free of trained values. Had we
adopted the more conventional integer pipeline with a learned
rescale~\cite{jacob}, a trained constant would necessarily appear in the
datapath and the clean separation between fabric and parameter image would be
lost. The restriction costs accuracy --- 0.07 points here --- and buys a
property that is easy to verify mechanically and hard to lose by accident.

The synthesis result reinforces this. We wrote 16 multiplications and got zero
multiplier cells on two independent targets. The architecture's job was not to
avoid multipliers directly but to arrange the arithmetic so that a standard
synthesis tool could remove them. That suggests a general design heuristic for
small quantized networks: expose constant operands to the tool and let it do
the strength reduction, rather than hand-coding shift-and-add.

The negative area result deserves to be carried forward rather than buried.
Over 1{,}024--50{,}176 bits an open-source mask ROM was consistently larger than
synthesized standard-cell storage, and the ratio flattened near $2.9\times$
rather than converging. We cannot attribute this: it may reflect the generator's
peripheral overhead, the library, or the size range, and distinguishing those
would require either a denser ROM generator or measurements well beyond the
sizes we could build. What the data does support is a caution --- the intuition
that a hard ROM is obviously denser than synthesized logic did not hold for any
size we could measure.

\section{Future Work}

The most valuable extensions, in rough order of leverage: importing models from
interchange formats such as ONNX rather than a single framework; supporting
arbitrary depth, which requires generalizing the two-phase control FSM and the
one-bit layer field in the memory interface; convolution lowering; and
place-and-route with timing analysis on both targets, which is the missing
evidence needed before any performance claim could be made. A trustworthy
physical-verification environment would also be required before the mask-ROM
path could be taken seriously.

\section{Limitations}

We are deliberate about what is not demonstrated.

\textbf{No physical signoff.} DRC and LVS results in our environment are not
trustworthy: OpenRAM's \emph{own} upstream reference ROM fails here with 830
DRC errors and an LVS mismatch. Consequently no result we produced is evidence
about our macros in either direction, and no macro is called DRC- or LVS-clean.
Physical \emph{generation} passes; physical \emph{signoff} is UNVERIFIED. These
are separate verdicts and we report them separately.

\textbf{No place-and-route or timing.} Neither target was placed, routed or
timed. No clock frequency is claimed; cycle counts are exact but any MHz figure
would be fabricated. The area comparison is between a hard-macro bounding box
and a standard-cell area sum, and is not a finished-block ratio.

\textbf{Scope.} MNIST only, one topology family, two dense layers, 4-bit
weights and uint8 activations. No convolution, no ONNX/TFLite ingestion. The
banking scheme is specific to the demonstrated ROM shape.

\textbf{Prior art.} This work explores a digital RTL interpretation of publicly
disclosed high-level Multiply-Select-Add ideas associated with public trade
coverage of commercial hardwired-inference silicon~\cite{taalas}. That source is
industry press, not peer-reviewed, and we grade it accordingly. No proprietary
source code, netlist, layout or transistor-level mask-ROM detail was used,
consulted or reproduced, and nothing here is claimed to be equivalent to that
hardware.

\section{Conclusion}

model2rtl demonstrates that a quantized MLP can be represented as a portable,
weight-index-driven RTL compute fabric whose functional behavior remains
bit-exact across behavioral simulation, FPGA-oriented synthesis and generic
logic synthesis. The enabling design decision is arithmetic, not architectural:
by fixing the requantization operator to a power-of-two shift and admitting no
multiplicative scale, the compute fabric becomes provably free of trained
values, and the model reduces to a parameter image behind a fixed interface.
The measured outcome --- zero multiplier cells after synthesis on both targets,
and zero functional mismatches on 500 images at every level of verification ---
suggests that for small quantized networks the interesting question is not how
to build multipliers efficiently, but how to arrange the arithmetic so that
synthesis can remove them.

The negative area result is worth carrying forward. Over the sizes we could
measure, an open-source mask ROM did not beat synthesized standard-cell
storage. Whether that reflects the generator, the library, or the size range is
the natural next question.

\section*{Reproducibility and AI Disclosure}

All RTL, parameter images, per-stage JSON reports and the full test suite are
part of the artifact. Every number in this paper was extracted programmatically
from those reports rather than transcribed. Quantities recorded by more than one
stage were cross-checked rather than reconciled. AI-assisted tools were used in
the development, verification and preparation of this work; all reported
measurements were produced by the named EDA tools and are reproducible from the
artifact.


\begin{thebibliography}{99}

\bibitem{jacob} B. Jacob, S. Kligys, B. Chen, M. Zhu, M. Tang, A. Howard,
H. Adam, and D. Kalenichenko, ``Quantization and training of neural networks for
efficient integer-arithmetic-only inference,'' in \emph{Proc. IEEE/CVF Conf.
Computer Vision and Pattern Recognition (CVPR)}, 2018, pp. 2704--2713.

\bibitem{deepcompression} S. Han, H. Mao, and W. J. Dally, ``Deep compression:
Compressing deep neural networks with pruning, trained quantization and Huffman
coding,'' in \emph{Proc. Int. Conf. Learning Representations (ICLR)}, 2016.

\bibitem{ucnn} K. Hegde, J. Yu, R. Agrawal, M. Yan, M. Pellauer, and
C. W. Fletcher, ``UCNN: Exploiting computational reuse in deep neural networks
via weight repetition,'' in \emph{Proc. 45th ACM/IEEE Int. Symp. Computer
Architecture (ISCA)}, 2018, pp. 674--687.

\bibitem{deepshift} M. Elhoushi, Z. Chen, F. Shafiq, Y. H. Tian, and J. Y. Li,
``DeepShift: Towards multiplication-less neural networks,'' arXiv:1905.13298,
2019.

\bibitem{lutnet} E. Wang, J. J. Davis, P. Y. K. Cheung, and G. A. Constantinides,
``LUTNet: Learning FPGA configurations for highly efficient neural network
inference,'' \emph{IEEE Trans. Computers}, vol. 69, no. 12, pp. 1795--1808,
2020.

\bibitem{logicnets} Y. Umuroglu, Y. Akhauri, N. J. Fraser, and M. Blott,
``LogicNets: Co-designed neural networks and circuits for extreme-throughput
applications,'' in \emph{Proc. 30th Int. Conf. Field-Programmable Logic and
Applications (FPL)}, 2020.

\bibitem{finn} Y. Umuroglu, N. J. Fraser, G. Gambardella, M. Blott, P. Leong,
M. Jahre, and K. Vissers, ``FINN: A framework for fast, scalable binarized
neural network inference,'' in \emph{Proc. ACM/SIGDA Int. Symp.
Field-Programmable Gate Arrays (FPGA)}, 2017, pp. 65--74.

\bibitem{hls4ml} J. Duarte \emph{et al.}, ``Fast inference of deep neural
networks in FPGAs for particle physics,'' \emph{Journal of Instrumentation},
vol. 13, no. 07, P07027, 2018.

\bibitem{openram} M. R. Guthaus, J. E. Stine, S. Ataei, B. Chen, B. Wu, and
M. Sarwar, ``OpenRAM: An open-source memory compiler,'' in \emph{Proc. 35th
Int. Conf. Computer-Aided Design (ICCAD)}, 2016.

\bibitem{yosys} C. Wolf, J. Glaser, and J. Kepler, ``Yosys --- a free Verilog
synthesis suite,'' in \emph{Proc. 21st Austrian Workshop on Microelectronics
(Austrochip)}, 2013.

\bibitem{taalas} T. P. Morgan, ``Taalas etches AI models onto transistors to
rocket boost inference,'' \emph{The Next Platform}, Feb. 19, 2026. [Industry
press; not peer-reviewed.]

\end{thebibliography}

\end{document}
